\PassOptionsToPackage{heading=true}{ctex}
\documentclass[lang=cn,chinesefont=nofont,11pt,a4paper,bibend=bibtex]{elegantpaper}
\newcommand{\reportkind}{设计文档}
% Dududa submission reports. Upstream ElegantPaper is kept unmodified in vendor/.
\usepackage{tabularx,longtable,array,makecell}
\usepackage{fancyhdr,needspace,placeins,float}
\usepackage{tikz}
\usetikzlibrary{arrows.meta,positioning,fit,backgrounds,calc}
\usepackage{fvextra}
\usepackage{microtype}
\geometry{a4paper,left=25mm,right=25mm,top=25mm,bottom=24mm,headheight=15pt,headsep=8mm,footskip=11mm}
\setCJKmainfont{FandolSong-Regular.otf}[BoldFont=FandolSong-Bold.otf,ItalicFont=FandolKai-Regular.otf]
\setCJKsansfont{FandolHei-Regular.otf}[BoldFont=FandolHei-Bold.otf]
\setCJKmonofont{FandolFang-Regular.otf}
\setmonofont{lmmono10-regular.otf}[Scale=0.9]
\definecolor{ink}{HTML}{202D3A}
\definecolor{accent}{HTML}{284B69}
\definecolor{muted}{HTML}{596773}
\definecolor{rulegray}{HTML}{BCC6CE}
\definecolor{panel}{HTML}{F2F5F7}
\color{ink}
\hypersetup{unicode=true,colorlinks=true,linkcolor=accent,urlcolor=accent,citecolor=accent,pdfauthor={嘟嘟哒项目},pdfcreator={LaTeX / ElegantPaper}}
\linespread{1.28}
\setlength{\parskip}{0.25em}
\setlength{\emergencystretch}{2em}
\clubpenalty=10000
\widowpenalty=10000
\setlength{\tabcolsep}{5pt}
\renewcommand{\arraystretch}{1.32}
\newcolumntype{L}[1]{>{\raggedright\arraybackslash}p{#1}}
\newcolumntype{Y}{>{\raggedright\arraybackslash}X}
\BeforeBeginEnvironment{tabularx}{\par\noindent}
\AfterEndEnvironment{tabularx}{\par\addvspace{0.5\baselineskip}}
\ctexset{
  section={format={\Large\sffamily\bfseries\color{accent}},
    aftername={\hspace{0.65em}},beforeskip={2.0ex plus .5ex},afterskip=1.1ex},
  subsection={format={\large\sffamily\bfseries},
    aftername={\hspace{0.6em}},beforeskip={1.7ex plus .4ex},afterskip=0.7ex},
  subsubsection={format={\normalsize\sffamily\bfseries},aftername={\hspace{0.6em}}}
}
\setcounter{tocdepth}{2}
\setcounter{secnumdepth}{3}
\setlist[itemize]{leftmargin=2em,itemsep=0.25em,topsep=0.35em}
\setlist[enumerate]{leftmargin=2em,itemsep=0.25em,topsep=0.35em}
\captionsetup{font=small,labelfont=bf,labelsep=quad}
\pagestyle{fancy}
\fancyhf{}
\fancyhead[L]{\small\sffamily 嘟嘟哒}
\fancyhead[R]{\small\sffamily\reportkind}
\fancyfoot[L]{\footnotesize\color{muted}2026 年 9 月 · 技术材料}
\fancyfoot[R]{\small\thepage}
\renewcommand{\headrulewidth}{0.35pt}
\renewcommand{\footrulewidth}{0pt}
\fancypagestyle{plain}{\fancyhf{}\fancyfoot[R]{\small\thepage}\renewcommand{\headrulewidth}{0pt}}
\newcommand{\code}[1]{{\small\expandafter\nolinkurl\expandafter{\detokenize{#1}}}}
\newcommand{\source}[1]{\par\smallskip{\footnotesize\color{muted}依据：#1}\par}
\newcommand{\reportdate}{2026 年 9 月 6 日}
\newcommand{\repolink}{\href{https://github.com/ly382965/dududa}{github.com/ly382965/dududa}}
\DefineVerbatimEnvironment{Code}{Verbatim}{fontsize=\footnotesize,breaklines=true,breakanywhere=true,frame=single,rulecolor=\color{rulegray},framesep=3mm,xleftmargin=1mm,xrightmargin=1mm}
\tikzset{box/.style={draw=rulegray,fill=panel,rounded corners=1pt,align=center,minimum height=11mm,font=\small,inner sep=5pt},flow/.style={-{Stealth[length=2mm]},draw=accent,thick},future/.style={box,dashed,fill=white},labeltext/.style={font=\footnotesize,color=muted,align=center}}
\newcommand{\reportcover}[3]{%
  \begin{titlepage}
  \thispagestyle{empty}\color{ink}
  \noindent{\small\sffamily\color{accent}DUDUDA\quad /\quad TECHNICAL DOCUMENTATION}\par
  \vspace{22mm}
  \noindent{\fontsize{34}{44}\selectfont\sffamily\bfseries 嘟嘟哒\par}
  \vspace{4mm}
  \noindent{\fontsize{24}{32}\selectfont\sffamily #1\par}
  \vspace{8mm}\noindent{\color{accent}\rule{\textwidth}{1pt}}\par
  \vspace{6mm}\noindent{\Large 群聊原生社交智能体\par}
  \vspace{3mm}\noindent{\large 面向校园生活、教学支持与群体协作\par}
  \vspace{15mm}
  \noindent\begin{minipage}{0.94\textwidth}\normalsize #2\end{minipage}
  \vfill
  \noindent\begin{tabularx}{\textwidth}{@{}L{27mm}Y@{}}
  提交材料 & #3\\
  文档日期 & \reportdate\\
  项目入口 & \repolink\\
  实现依据 & 本地源码与已存验证记录；能力状态详见正文\\
  \end{tabularx}
  \vspace{12mm}
  \end{titlepage}
}

\hypersetup{pdftitle={嘟嘟哒 设计文档},pdfsubject={群聊智能体的系统架构、功能实现与技术验证}}
\begin{document}
\hypersetup{pageanchor=false}
\reportcover{设计文档}{本报告说明嘟嘟哒的使用场景、系统架构、核心模块、模型调用、工程实现与验证结果。内容以当前主仓实现为依据，整合本地设计、研究、原型、评测及运维资料，并明确区分运行能力、独立组件与后续设计。}{第 1 项 · 技术设计报告}
\pagenumbering{Roman}
\hypersetup{pageanchor=true}
\section*{摘要}
\phantomsection
\addcontentsline{toc}{section}{摘要}
校园群聊同时承载课程咨询、活动协调、通知转发与日常交流。面对多人交错发言、短句省略、临时更正及分散的数据源，智能体需要连续完成情境理解、参与决策、资料查询和回答组织。嘟嘟哒 以 QQ 群为主要交互场景，将这些职责拆分为上下文与感知、Social Engine、模型路由、Capability/MCP、回答与 Persona、输出适配等模块，并通过统一 Python Runtime 编排。

系统经 NapCat/OneBot v11 与 AstrBot 接入 QQ，使用 Vue 3/Node.js 控制台管理群服务、模型配置与运行记录。业务层已连接评课、第二课堂、教务、培养方案、校园通知与校车六类校园服务。模型调用采用 DeepSeek 的 OpenAI 兼容 Chat Completions 接口，模型档位、推理强度和可见回答篇幅分别配置。管理员授权后，系统可以依据连续讨论和当前请求启用匹配的只读查询插件。

本报告以自适应机制为技术主线：根据连续讨论与当前请求选择已授权插件；根据任务复杂度、置信度和预算选择模型路由；根据群聊情境调整参与方式与表达；并以 WebUI 工作台呈现实际选择，支持按群配置、预览和调整。Memory v2 的作用域检索、中文 BM25 和生命周期组件，以及主动调度、离线 Bandit 等研究成果一并说明，但不将其等同于已完成的长期群聊体验。

验证材料包含 100 条固定消息的离线流程、真实 MCP 原子查询、真实模型预览及模拟入站记录。报告分别呈现测试输入、实际执行范围、结果与限制，保留早期广覆盖预览中发现的部分通过和失败项，为评审者提供可核查的完成度说明。

\noindent\textbf{关键词：}群聊智能体；校园服务；语义感知；模型路由；MCP；可观察运行时
\vfill
\noindent\textbf{阅读说明。}第 1—4 节介绍目标、场景与架构；第 5—9 节解释技术机制；第\ref{sec:webui}节独立介绍 WebUI 设计；第\ref{sec:deployment}—\ref{sec:evaluation}节说明部署与验证；第\ref{sec:synthesis}节及附录汇总本地资料和演进边界。文中方括号 S 编号对应附录的本地来源索引。
\clearpage
\begingroup\small\linespread{1.05}\selectfont\setlength{\parskip}{0pt}
\tableofcontents
\endgroup\clearpage
\pagenumbering{arabic}
\section{作品目标与设计思路}
\subsection{作品定位}
嘟嘟哒 是面向持续群体交流的社交智能体。其产品定位为“生活在 QQ 群里的 AI 群友”，技术目标是在同一条运行链中完成理解群聊情境、判断参与方式、使用外部能力、形成回答并可靠投递。校园生活与教学支持是首批应用场景，系统的接口划分同时支持向其他群体服务扩展。

群聊具有与单人问答不同的条件：当前消息可能只是一句追问，指向另一人的发言；同一时间可能出现多个话题；群成员之间的更正具有先后关系；一次有用的回答也可能在错误的时机造成打扰。因此，作品把“是否回应、回应谁、需要什么资料、如何表达”分别建模，使技术实现能够围绕真实交互问题调整。

\subsection{需要解决的问题}
\begin{enumerate}
\item \textbf{信息分散。}课程评论、活动、教务安排和通知分布于不同系统，用户需要反复切换页面，并自行核对课程身份、时间范围和资料版本。
\item \textbf{上下文不完整。}省略、引用、别称和更正使当前短句难以独立理解；仅拼接历史又可能把旧安排当成现行事实。
\item \textbf{参与缺乏节奏。}逐条回复会干扰讨论，单纯等待指令则难以利用连续讨论中已经形成的需求。
\item \textbf{难以解释运行过程。}模型、查询服务、配置和消息平台分别出错时，管理员需要知道具体失败环节，以及保存的配置是否实际生效。
\end{enumerate}

\subsection{设计原则与工程选择}
\textbf{明确分工。}输入适配器确认平台事件，感知模块解释语义，社交策略选择行动，能力层获取资料，回答层组织内容。身份、群范围、执行权限和输出目标由程序维护，模型承担语义提取与文本生成。

\textbf{以业务能力连接异构服务。}对智能体暴露的是课程查询、活动搜索、通知读取等业务单元。具体 MCP 工具名、参数投影和服务会话留在适配层。这样既能复用校验与错误处理，也便于保留各数据源的差异。

\textbf{按任务分配计算。}模型档位、推理强度与输出篇幅独立配置，避免将“详细回答”直接解释为“必须换更大的模型”。感知与回答分别选择模型角色和预算。

\textbf{让运行可观察。}QQ 与 Web 预览共用 Runtime；控制台展示实际运行模型、能力、耗时及结果。开发验证使用同一套编排和可替换适配器，减少独立演示实现带来的偏差。

\subsection{本次报告的实现口径}
报告以 2026 年 9 月 6 日本地主仓为对象。开始核对时主仓提交为 \code{2b52593}，工作区还包含录制与查询相关的在途改动；后续章节同时引用明确日期的已存验证材料。提交号用于定位源码，不表示所有历史测试都运行于该提交，也不表示本次重新发布了程序。

本文将实现分为三类：\textbf{已进入当前运行组合}，如校园查询、模型路由和按群自动参与；\textbf{已完成独立组件或离线验证}，如 Memory v2 和 Bandit 估计器；\textbf{后续设计}，如通用多步 Skill、长期群友印象与自动学习。各项验证的运行模式在第\ref{sec:evaluation}节单独列出。
\source{[S1, S2, S3, S10]。历史原型、研究报告与当前主仓具有不同适用范围。}

\section{应用场景与需求分析}
\subsection{用户与使用方式}
主要使用者包括课程学习群中的学生、活动组织者、讨论组成员，以及维护智能体的群管理员。群成员通过自然语言提问；明确提及 Bot 时进入直接请求流程，开启自动参与后，连续群聊也可以构成参与机会。管理员负责授予群服务权限、配置参与节奏，并在控制台检查运行情况。

\begin{table}[htbp]\centering
\caption{校园与群体协作场景}
\begin{tabularx}{\textwidth}{@{}L{22mm}Y Y@{}}\toprule
场景 & 典型问题 & 系统应完成的工作\\\midrule
教学支持 & “数学分析 B1 有哪些教师，评论分别怎样？” & 提取课程线索，查询课程与教师，区分评论样本和客观开课信息，保留来源。\\
活动参与 & “二课最近有哪些学术活动，最多三项？” & 识别类别、时间和数量限制，返回活动名称、时间、状态与详情入口。\\
学业规划 & “2026 级计算机普通主修需要多少学分？” & 指定年级、专业和方案类型；返回相应版本资料，说明研究快照口径。\\
校园生活 & “工作日东区到西区上午有哪些校车？” & 解析方向、日期类型、时段；区分始发与途经时间，注明时刻表版本。\\
服务与管理 & “最近三条校园通知是什么？” & 搜索通知，呈现标题、发布日期和原始链接；必要时查询截止日期。\\
群体协作 & “周五改成周六了，最终什么时候开会？” & 使用同群近期历史识别更正关系，并针对当前问题给出明确结论。\\\bottomrule
\end{tabularx}
\end{table}

科研讨论群可以使用一般问答、分工草案和讨论整理能力；当前材料尚未提供专业科研检索、实验管理或科研成果提升的专项验证，因此作品价值主要落在教学支持、校园生活和群体协作。

\subsection{功能需求}
一次校园查询应能从自然语言定位业务服务，绑定必要参数，读取有来源的资料并形成可理解的回答。信息不足时需要澄清；远端服务失败时应显示失败原因或可用的降级说明，不能把未取得的资料补写为事实。

一次群聊参与应保留发言者、引用关系和当前问题，受到本群配置约束。自动启用只读插件需要已有管理员授权，不能因发现一个新话题就安装任意代码或扩大服务范围。

一次管理操作应让用户看到配置目标和实际状态。QQ 连接、账号登录、Agent 可用、模型配置和 MCP 服务连接应分别表达，因为其中任一项可独立失效。

\subsection{质量需求与可检查目标}
\begin{tabularx}{\textwidth}{@{}L{24mm}Y@{}}\toprule
质量维度 & 可检查目标\\\midrule
语义与事实 & 更正优先于旧安排；课程与教师不混淆；结果中的年份、来源和数量限制保持一致。\\
参与节奏 & 自动参与受概率、冷却和小时上限约束；被关闭的群或能力不能绕过配置执行。\\
可维护性 & 核心包不依赖 AstrBot 的事件对象；新增服务主要修改能力定义、映射和 Provider。\\
可诊断性 & 每次运行能区分模型失败、工具失败、未参与和已生成但未投递等状态。\\
可部署性 & 提供依赖锁文件、容器编排和配置示例；离线流程可在无模型 Key、无 QQ 登录时运行。\\\bottomrule
\end{tabularx}

这些目标用于组织实现与验证，尚未形成真实群长期使用的服务等级承诺。当前没有足够样本支持总体准确率、长期打扰率或 P95 延迟指标。
\source{[S1, S8, S10, S11]。场景示例为说明输入，不是未经执行的效果承诺。}

\section{总体架构与运行流程}
\subsection{系统边界与分层}
系统采用单仓组织智能体核心、QQ 插件、Web 控制台、MCP 服务和部署工具。Python 核心通过 Domain 数据类型与 Ports 接口表达消息、策略、模型、能力和输出，AstrBot 组合根负责装配具体实现。外部模型平台、校园认证系统和校园数据源提供服务；公共域名、统一登录与站点网关属于部署环境，不是智能体核心组件。

\begin{figure}[H]\centering
\begin{tikzpicture}[x=1cm,y=1cm]
\node[box,text width=6.1cm] (qq) at (0,0) {QQ 群成员\\NapCat / OneBot v11};
\node[box,text width=6.1cm] (web) at (7.1,0) {管理员 Web 工作台\\Vue 3 / Node.js Gateway};
\node[box,text width=6.1cm] (ad) at (0,-1.65) {AstrBot 输入与输出适配\\消息身份、引用、群历史、投递};
\node[box,text width=6.1cm] (pr) at (7.1,-1.65) {Runtime 预览与配置接口\\服务端同群历史、运行记录};
\node[box,text width=13.2cm,minimum height=2.1cm] (rt) at (3.55,-3.9) {\textbf{Dududa Agent Runtime}\\[2mm]上下文 $\rightarrow$ 感知 $\rightarrow$ 社交策略 $\rightarrow$ 模型路由\\[1mm]能力计划与执行 $\rightarrow$ 回答组织 $\rightarrow$ Persona $\rightarrow$ 输出请求};
\node[box,text width=3.75cm,minimum height=1.35cm] (llm) at (-1.05,-6.4) {模型适配器\\DeepSeek\\Chat Completions};
\node[box,text width=3.75cm,minimum height=1.35cm] (mcp) at (3.55,-6.4) {Unified MCP / 本地能力\\五类 MCP 校园服务\\校车时刻表};
\node[box,text width=3.75cm,minimum height=1.35cm] (state) at (8.15,-6.4) {策略与状态\\Persona / 群配置\\SQLite / JSON};
\node[future,text width=13.2cm] (future) at (3.55,-8.15) {独立模块与后续接入：Memory v2、主动日报、离线 Bandit、多步 Skill};
\draw[flow] (qq)--(ad);\draw[flow] (web)--(pr);
\draw[flow] (ad.south)--(ad.south |- rt.north);
\draw[flow] (pr.south)--(pr.south |- rt.north);
\draw[flow] (rt.south -| llm.north)--(llm.north);
\draw[flow] (rt)--(mcp);
\draw[flow] (rt.south -| state.north)--(state.north);
\draw[flow,dashed] (future.east) -- ++(0.4,0) |- (rt.east);
\end{tikzpicture}

\caption{嘟嘟哒的系统组成。实线为当前交互或调用关系；虚线模块为独立实现或后续接入方向。}
\label{fig:architecture}
\end{figure}

\begin{longtable}{@{}L{24mm}L{45mm}L{80mm}@{}}
\caption{架构分层与责任}\label{tab:layers}\\\toprule
层次 & 主要组件 & 责任与输入输出\\\midrule\endfirsthead
\toprule 层次 & 主要组件 & 责任与输入输出\\\midrule\endhead
交互与适配 & QQ、NapCat、AstrBot；Vue 3、Node.js & 将平台消息变为统一消息，将结果变为 QQ 消息或 Web 预览；控制台读取配置和运行记录。\\
领域与端口 & Domain、Ports & 表达 Actor、ConversationScope、MessageEnvelope、调用预算和版本化结果；不携带框架事件或密钥。\\
智能体编排 & Runtime、Perception、Social、Router & 管理阶段流转，构建上下文，识别意图，决定静默、澄清、回答或使用工具，选择计算配置。\\
业务能力 & Registry、Planner、Executor、Validator & 筛选本群可用能力，绑定输入，执行有界查询，将结果整理为 Observation。\\
内容与表达 & DirectChat、Composer、Persona、Final Validator & 使用问题或可信工具资料生成正文；组织结构、来源与篇幅；完成最终格式校验。\\
基础设施 & 模型适配器、Unified MCP、SQLite/JSON、Compose & 提供外部访问、服务会话、配置和持久状态；隔离第三方依赖及运行数据。\\\bottomrule
\end{longtable}

\subsection{核心组件与依赖方向}
组合根位于 \code{apps/astrbot-plugins/astrbot_plugin_dududa_core/}，其中 \code{composition.py} 创建规则与模型混合感知、确定性社交策略、静态模型路由、单步能力运行时、Persona Registry 及输出请求构建器。

核心 Runtime 的类名仍为 \code{OfflineRuntimeOrchestrator}，沿用了早期离线实现名称。当前正式组合已经为它注入真实模型和 MCP 适配器，因此类名中的 Offline 不代表正式运行使用假模型。测试则可以替换这些接口而保留相同编排。

\begin{center}
\begin{tikzpicture}[node distance=7mm]
\node[box,text width=133mm] (a) {应用组合根与平台适配器\quad AstrBot / Web / Model / MCP};
\node[box,text width=133mm,below=of a] (b) {应用服务与 Runtime\quad 上下文、感知、社交、路由、能力、回答};
\node[box,text width=133mm,below=of b] (c) {Domain 与 Ports\quad 统一数据模型、接口、调用与结果契约};
\draw[flow] (a)--(b);\draw[flow] (b)--(c);
\end{tikzpicture}
\end{center}
核心依赖抽象端口，应用组合根提供具体实现。新增平台需要适配消息身份、引用和投递回执；新增模型需要满足统一请求与结果协议；新增业务则实现能力定义和 Provider。这种拆分支持独立验证，但也要求维护类型、组合根与接口实现之间的一致性。

\subsection{关键数据对象}
\begin{tabularx}{\textwidth}{@{}L{38mm}Y@{}}\toprule
对象 & 业务含义\\\midrule
MessageEnvelope & 本轮平台消息、正文、时间、提及、引用及附件描述；作为运行输入。\\
Actor / ConversationScope & 当前行为主体及平台、机器人、会话、群的范围；决定可读取和可回复的上下文。\\
PerceptionResult & 意图、实体、引用关系、能力类别及置信信息；提供后续决策所需的结构化语义。\\
SocialDecision / ResponsePlan & 参与动作与理由、回答目标、模型档位和表达预算；区分行为选择与文本生成。\\
ToolPlan / Observation & 待执行的业务能力与参数，以及校验后的工具事实、来源和结果状态。\\
Draft / FinalResponse & 回答结构与最终正文；Persona 和格式校验在这一层完成。\\
DeliveryRequest / Receipt & 期望投递的目标、分片和内容，以及平台确认的实际结果。\\\bottomrule
\end{tabularx}

\subsection{一次查询的完整路径}
以“评课社区里数学分析 B1 的评价怎样”为例，系统首先从事件取得可信身份和同群上下文。规则识别明确的评课来源，快速档模型提取课程实体和意图；社交策略允许回应，能力目录筛选本群可用的评课能力；参数绑定器形成一次业务查询。MCP 返回课程、教师、评论和链接后，回答模型按事实组织文字，Composer 和 Persona 处理表达，最终由对应输出适配器完成呈现。

\begin{figure}[H]\centering
\begin{tikzpicture}[x=1cm,y=1cm]
\foreach \x/\name/\id in {0/入口适配/a,3.5/Runtime/b,7/模型适配/c,10.5/能力与 MCP/d}{
  \node[box,text width=2.6cm] (\id) at (\x,0) {\name};
  \draw[dashed,rulegray] (\x,-0.6)--(\x,-7.25);
}
\draw[flow] (0,-1)--node[above,font=\footnotesize]{统一消息与同群上下文}(3.5,-1);
\draw[flow] (3.5,-2)--node[above,font=\footnotesize]{结构化感知请求}(7,-2);
\draw[flow] (7,-2.7)--node[above,font=\footnotesize]{意图、实体与置信信息}(3.5,-2.7);
\draw[flow] (3.5,-3.6)--node[above,font=\footnotesize]{资格筛选、参数绑定、单步执行}(10.5,-3.6);
\draw[flow] (10.5,-4.4)--node[above,font=\footnotesize]{Observation：事实、身份、来源、状态}(3.5,-4.4);
\draw[flow] (3.5,-5.25)--node[above,font=\footnotesize]{资料与回答约束}(7,-5.25);
\draw[flow] (7,-5.95)--node[above,font=\footnotesize]{生成正文与用量}(3.5,-5.95);
\draw[flow] (3.5,-6.8)--node[above,font=\footnotesize]{格式校验后输出}(0,-6.8);
\end{tikzpicture}

\caption{校园查询的请求时序。模型承担感知与回答，工具调用由程序显式执行。}
\label{fig:sequence}
\end{figure}

一次典型查询由感知、业务能力执行和回答生成构成。业务能力内部可能读取多个数据条目，但当前 Planner 的计划仍只包含一个业务步骤；这与通用智能体跨服务自主规划多步任务不同。多项独立需求超出本轮范围时，当前实现会澄清或引导拆分。

\subsection{阶段状态与完成条件}
\begin{figure}[H]\centering
\begin{tikzpicture}[node distance=7mm and 7mm]
\node[box,text width=3.8cm] (r) {接收与预处理\\received / preprocessed};
\node[box,text width=3.8cm,right=of r] (c) {上下文与感知\\context\_ready / perceived};
\node[box,text width=3.8cm,right=of c] (d) {参与决策\\decided};
\node[box,text width=3.8cm,below=of d] (t) {计划与执行\\tools\_planned / executed};
\node[box,text width=3.8cm,left=of t] (v) {结果与内容\\validated / composed};
\node[box,text width=3.8cm,left=of v] (p) {表达与输出准备\\rendered / ready\_to\_emit};
\node[box,text width=13cm,below=of v] (end) {输出回执 $\rightarrow$ 完成\quad delivery\_acknowledged $\rightarrow$ completed};
\draw[flow] (r)--(c);\draw[flow](c)--(d);\draw[flow](d)--(t);\draw[flow](t)--(v);\draw[flow](v)--(p);
\draw[flow](p.south)--(p.south |- end.north);
\end{tikzpicture}

\caption{主要运行阶段。直接回答跳过工具阶段；错误和延期可从相应阶段终止。}
\end{figure}

运行状态记录从收到消息到生成最终内容的进展，输出请求与实际发送分开确认。\code{ready_to_emit} 只表示内容已准备好，\code{delivery_acknowledged} 表示收到输出适配器回执。Web 预览把结果交给预览适配器，不能据此推断 QQ 已投递。状态类型中保留 Memory 阶段，但当前生产组合未装配长期记忆读写。
\source{[S2, S3, S4]；源码：\nolinkurl{runtime/state.py}、\nolinkurl{runtime/orchestrator.py}、\nolinkurl{composition.py}。}

\section{功能模块与业务能力}
\subsection{六类校园服务}
当前运行组合将五类 MCP 校园服务与校车本地 Provider 接入统一能力层。这里的“六类”按面向用户的业务领域统计，不等于 MCP 工具数量，也不包括控制台中仅登记的可选服务。

\subsubsection{评课社区}
评课能力支持课程与教师搜索、课程详情、评论、排名及统计。高层查询入口 \code{icourse.public-query.v2} 将业务操作与查询条件映射至 iCourse 服务。回答需要区分同名课程、不同教师、评论条数与站点统计；评论中的主观体验应表述为评论观点，不能转化为对课程或教师的确定性评价。

在选课场景中，多位群友可以分别提出讲解细致程度、证明训练和作业负担等偏好。工具提供课程与评论材料，回答模型负责按这些维度组织建议。若评论没有覆盖某一维度，应说明材料不足；学期是否实际开课属于教务信息，不能仅靠评课社区推断。

\subsubsection{第二课堂}
二课服务提供活动搜索、详情、类别字典及连接状态。参数绑定处理活动关键词、类别、时间和数量，返回活动名称、时间、报名状态与资料范围。运行环境通过部署端 CAS 凭据访问服务；群聊问答面向活动信息，个人学时、已报名记录、签到及代报名不属于本次智能体查询范围。

自然语言中“最近”“公开发布”“学术活动”等条件需要分别落实为查询或结果过滤。修复记录表明，类别未绑定、活动详情参数错误或返回院系树过大，都会导致错误内容或过高的上下文消耗。项目针对这些具体问题收敛了参数和返回结构。

\subsubsection{教务公开信息}
教务服务支持学期列表、全校开课、考试和教学日历。课程名称与教师来自开课记录，学期和日期字段随结果保留。检索失败、空结果与没有考试安排是不同情况，回答应按实际返回区分。当前接口不提供学生个人成绩查询，也不替代校方教务操作。

\subsubsection{培养方案研究资料}
培养方案服务按年级、专业、课程代码和方案类型查询，并支持方案与课程比较。数据来自公开资料整理的非官方研究快照，当前来源覆盖 2015—2026 级，快照日期为 2026 年 8 月 2 日。课程身份使用精确课程代码，防止同名课程在不同方案中被误合并。

该服务适合查阅和讨论培养要求。报告中的学分结果必须绑定年级、专业与方案类型；不能把快照当成实时教务数据库，或把一次查询当成个人毕业资格审核。

\subsubsection{校园通知}
NotifAI 服务提供通知搜索、详情、日历、截止日期、来源与分类字典及统计。不同操作拥有独立能力定义，查询内容与查询字典应分开。回答保留标题、发布日期、截止时间和来源链接，在时间字段缺失时明确说明。

“最近三条通知”属于通知正文检索，不能误选“列出通知来源”。相关录制回归针对这类口语到工具操作的误选作了修复，说明工具成功返回并不自动意味着业务选择正确。

\subsubsection{校车查询}
校车采用本地版本化时刻表，通过 \code{ustc.shuttle.public-query.v1} 进入 Runtime。插件不注册独立自然语言消息处理器，意图识别、查询、总结和发送仍由核心负责。查询支持校区、路线、日期类型和时段，返回相应班次。

当前资料中东区、西区、南区及高新校区版本日期为 2026 年 8 月 30 日，太湖路园区为 8 月 27 日。实现区分始发站与经过站时间，并保留班次备注；“没有公交车辆运行”不等于该班次取消。时刻表通过维护数据文件更新，属于静态资料查询。

\subsection{近期对话、协作与表达}
直接交流可以用于解释问题、分析方案、整理多人分工和回答上下文追问。近期讨论总结从可获取的同群历史构建，控制台显示读取条数和覆盖范围；若只有最近若干条消息，就不应称为“全天完整摘要”。分工草案需要将人员、任务和产出对应起来，不能只列抽象角色。

Persona 资产使同一智能体在技术讨论与日常交流中保持连贯语气。精确输出要求优先于人格表达：要求只回复指定文字时不得附加语气词，要求 JSON 时不得插入解释或表情。长回答由输出层按平台能力处理。

\subsection{群插件与可选服务}
\begin{longtable}{@{}L{31mm}L{51mm}L{67mm}@{}}
\caption{扩展资产与当前使用范围}\\\toprule
组件 & 已有功能 & 本次状态与限制\\\midrule\endfirsthead
\toprule 组件 & 已有功能 & 本次状态与限制\\\midrule\endhead
Emoji Kitchen & 显式表情组合，返回合成图像 & 已有真实服务生成素材；与语言模型生成图片分开。\\
复读插件 & 根据重复消息规则响应 & 独立、按群配置；录制固定样例验证处理器，不代表长期自然群表现。\\
Arcaea 扩展 & 曲目信息、本地谱面、绑定及 B50 相关资产 & 曲目与谱面有已生成素材；真实 B50 查分本次保持关闭；本地 B50 Provider 未接入生产组合。\\
Sub2API Readonly & 显式服务统计命令 & 录制使用真实处理器处理固定统计样例，不能将其当成实时业务统计。\\
社交扩展插件 & 生日、睡眠记录、投票等显式命令与状态组件 & 默认关闭，独立于 Core；不自动发送生日祝福或情绪安慰。\\
五个可选 MCP & 本地推荐、专业设置、学校通知、学院通知、图书馆开放时间 & 已有服务与目录接入、连接检测记录；未配置相应 Capability mapping 时，不进入当前 Planner 候选。\\
旧兼容组件 & Target Talk、Reply Polish、旧记忆等源码与迁移材料 & 不据源码存在推断当前启用；当前主链由 Core 负责入站与输出。\\\bottomrule
\end{longtable}

将独立插件、可发现服务和可规划能力分开描述，可以明确回答“功能已写好”“服务能连接”“用户可以通过自然语言调用”这三个不同问题。扩展资产为后续接入提供基础，本次成果以已经贯通的业务路径为主。
\source{[S5, S6, S7, S10]；各校园服务与插件 README，原子查询和素材记录。}

\section{上下文感知与群聊参与机制}
\subsection{上下文范围与消息解释}
消息适配保留消息标识、平台、机器人、会话、群、发言者、时间、提及与引用关系。Web 预览从服务端获取同账号、同群的近期历史；主动参与执行器读取当前群窗口。Runtime 继续按预算整理这些资料，区分本轮问题、被引用内容和历史发言。

本轮指令决定要完成的任务，历史发言解释指代和讨论背景。例如，旧消息提出周五开会，后续更正为周六，本轮追问“最终时间”，应依据有效更正形成答案。来源于其他群的引用文本可以作为本轮提供的资料，但不能成为读取其他群完整历史的授权。

\subsection{规则感知与模型感知合并}
规则适合识别明确提及、已知业务词和输入格式；模型适合处理课程别称、口语、省略与开放式需求。\code{HybridPerceptionEngine} 将两者组织为结构化感知结果，模型输出经过本地 Schema 校验后才能参与决策。当前感知并非向后续层直接透传任意自然语言推断。

结构化结果用于表达意图、实体、引用、偏好能力类别及置信信息。明确的“评课社区”提示来源，课程实体提示查询对象；“最多三项”属于结果数量约束。模型只提供语义候选，最终可用能力、执行预算和群权限仍由程序检查。

复杂度评估与社交决策分别使用感知信号。当前组合采用确定性策略，不为每一个阶段额外调用一个大模型。高复杂度档位选择需要相应信号和置信条件，低置信、指代不明或任务缺项可以进入澄清流程。

\subsection{Social Engine 的动作选择}
Social Engine 将对内容的理解转化为参与动作。当前核心策略包括直接回应、调用业务能力、澄清和保持静默等行为路径；输入权限、群模式与当前请求条件影响选择。它与外围自动参与执行器职责不同：前者决定本轮怎么处理，后者先判断普通群消息是否值得启动这一轮。

显式提及和自动参与使用不同入口。用户明确提出任务时，系统进入请求处理；自动参与则在目标群已开启的前提下检查事件类型、概率、冷却、每小时上限及是否有进行中的运行。随机抽样只影响参与机会，不应改变事实或服务权限。

\begin{table}[htbp]\centering
\caption{自动参与的上下文窗口配置（当前代码值）}
\begin{tabularx}{\textwidth}{@{}L{32mm}L{28mm}L{30mm}Y@{}}\toprule
窗口 & 最大消息数 & 字节预算 & 适用含义\\\midrule
compact & 12 & 3,500 & 紧凑讨论片段\\
standard & 30 & 6,500 & 一般群聊上下文\\
extended & 100 & 7,500 & 更多发言记录，仍受总预算约束\\\bottomrule
\end{tabularx}
\end{table}

窗口越大不等于可完整读入更多长消息，也不等于全天历史。当前控制器维护发送时间和运行占用状态；冷却可以抑制短时连续触发，但不能独立证明跨冷却窗口的建议不会语义重复。这一问题需要专门的长期群聊评估。

\subsection{受授权的能力自适应启用}
当前自适应选择器是可解释的规则实现。它只在管理员允许的只读插件集合中匹配需求；并要求足够的连续讨论、多位发言者和当前明确请求。核对时的条件包括：
\begin{enumerate}
\item 历史至少 8 条，解析得到至少 3 个发言者标签；
\item 最后一条消息包含查询、建议、推荐或疑问等请求信号；
\item 最近 20 条中至少 3 条与同一候选主题相关，最近 4 条仍出现该主题；
\item 候选插件已经获得本群自适应启用授权。
\end{enumerate}

规则满足后，策略存储记录能力启用和原因，当前请求进入正式 Runtime，控制台读取实际生效状态。该机制能够在讨论逐渐形成需求时开放已有能力；它不下载新插件，不修改程序，也不是在线训练或自主生成 Skill。

\begin{figure}[H]\centering
\begin{tikzpicture}[node distance=6mm]
\node[box,text width=13cm] (a) {群授权 + 连续讨论窗口 + 当前请求};
\node[box,text width=13cm,below=of a] (b) {规则匹配：发言者标签、主题密度、近期相关性、允许的只读插件};
\node[box,text width=13cm,below=of b] (c) {记录启用与原因 $\rightarrow$ Runtime 查询 $\rightarrow$ Web 状态同步};
\draw[flow](a)--(b);\draw[flow](b)--(c);
\end{tikzpicture}
\caption{当前自适应启用机制。启用范围来自管理员配置，语义回答由后续 Runtime 完成。}
\end{figure}

四人 20 条固定讨论证明了这一链路的代表场景：前 19 条评课关闭，第 20 条形成请求后启用评课。20 条是该场景的设置，不是算法对所有群的固定阈值；更多群聊类型、话题切换和误触发仍需单独评估。
\source{[S3, S4, S10]；\nolinkurl{adapters/proactive_talk.py}、\nolinkurl{adapters/adaptive_plugins.py} 与 \nolinkurl{perception/social.py}。}

\section{模型选择与 API 调用方式}
\subsection{模型档位、推理强度与回答篇幅}
系统使用三个独立维度描述计算与表达。Tier 选择模型配置；Reasoning 决定推理投入；AnswerProfile 约束用户可见的正文长度、段落数及投递形式。当前三档的界面名称为 Luna、Terra、Sol，内部枚举为 haiku、sonnet、opus。这些是程序别名，不能据此认定调用的是 Anthropic 模型。

\begin{table}[htbp]\centering
\caption{项目记录的常规 DeepSeek 配置}
\begin{tabularx}{\textwidth}{@{}L{22mm}L{22mm}L{43mm}L{17mm}Y@{}}\toprule
界面档位 & 内部 Tier & 实际模型 ID & 推理值 & 主要用途\\\midrule
Luna 快速 & haiku & \code{deepseek-v4-flash} & low & 感知、简短交流\\
Terra 标准 & sonnet & \code{deepseek-v4-flash} & high & 校园查询、一般分析\\
Sol 深度 & opus & \code{deepseek-v4-pro} & max & 复杂比较、论证\\\bottomrule
\end{tabularx}
\end{table}

2026 年 9 月 6 日录制查询环境保留快速档的轻量推理，标准档和深度档关闭额外 thinking，直接组织已取得的工具资料。因此，表中常规推理配置不能直接用于解释该批样本耗时；第\ref{sec:evaluation}节保留这一差异。真实问答样本主要覆盖 Flash，已有配置并不等于三档模型完成了公平的性能对比。

\subsection{任务驱动的自适应路由}
感知角色使用快速档模型，回答角色依据任务信号、群偏好和可用档位选择模型。静态路由器在调用前检查角色、能力、数据条件、输出容量和预算，形成可记录的路由结果；管理员也可以固定档位。当前路由没有调用在线 Bandit 学习器。

这里的自适应体现为同一套策略针对不同任务作出不同选择。当前档位策略使用复杂度信号及置信度：深度推理、独立核验、多约束综合等信号可支持选择较高档位；信号不足时使用默认或较低档位，并检查预算和候选可用性。模型档位与选择原因进入运行结果，便于管理员理解一次请求为什么采用相应计算配置。

推理额度与正文额度分别约束。如果回答只需一句话，模型仍可能需要充分推理；若用户需要详细步骤，也不必仅因篇幅较长就提升任务复杂度。当前适配器将内部 \code{light/deep/maximum} 映射到 DeepSeek 的 \code{low/high/max}。无额外推理时发送 \code{thinking: disabled}。

一次典型工具问答的等待时间可分解为
\[
T=T_{\mathrm{context}}+T_{\mathrm{perception}}+T_{\mathrm{tool}}+T_{\mathrm{answer}}+T_{\mathrm{output}}.
\]
这只是运行阶段的拆分，不是测得的耗时模型。现有样本记录了总耗时，尚不足以量化各阶段的平均贡献；后续优化应先区分工具等待与模型生成，再决定缓存或调用合并措施。

\subsection{调用链与请求示例}
\begin{center}\small
Runtime ModelRequest $\rightarrow$ AstrBotModelProviderAdapter\\
$\rightarrow$ AstrBot Provider.text\_chat $\rightarrow$ DeepSeek Chat Completions
\end{center}
服务基础地址为 \url{https://api.deepseek.com}，调用采用 OpenAI 兼容的 Chat Completions 协议。以下是与当前标准档录制模式对应的等价 HTTP 请求，用于说明字段；真实问题、系统提示和工具资料由 Runtime 构建。

\begin{Code}
POST https://api.deepseek.com/chat/completions
Authorization: Bearer <DEPLOYMENT_API_KEY>
Content-Type: application/json

{
  "model": "deepseek-v4-flash",
  "messages": [
    {"role": "system", "content": "Answer using the supplied evidence."},
    {"role": "user", "content": "<question and validated tool observations>"}
  ],
  "thinking": {"type": "disabled"},
  "max_tokens": 8192
}
\end{Code}

常规启用推理的请求改用 \code{"thinking": {"type": "enabled"}}，并添加 \code{"reasoning_effort": "high"} 等配置。DeepSeek 官方文档说明了这两个控制字段；如直接使用 OpenAI SDK，扩展字段 thinking 应通过 SDK 的 \code{extra_body} 传递。[W3]

项目本身通过 AstrBot \code{text_chat} 传入 \code{prompt}、\code{system_prompt}、模型 ID、\code{max_tokens} 和推理参数，由宿主 Provider 转换为 HTTP 请求。当前核心并不依赖模型原生 Function Calling 自动执行校园工具：语义提取完成后，由 Capability Runtime 调用 MCP，再将校验结果送入回答模型。

\subsection{响应解析与错误处理}
适配器读取 \code{completion_text}、token 用量及终止信息，转换为统一 ProviderResponse。感知结果须满足 JSON Schema；回答请求接收正文，再经过内容结构、篇幅和格式处理。推理正文不作为用户答案，也不需要作为普通群聊历史展示。

请求受本轮总预算和下游超时约束。认证失败、限流、超时、输出无效和用户取消应成为可区分的错误，控制台以公开错误码和说明展示。重试或切换端点只能发生在已配置策略和预算内，不能无限重复调用；网页返回成功状态也不能代替有效正文检查。

\subsection{Key 池与配置生效}
Web 提供三档 Key 池，维护 Provider、基础地址、协议、模型及凭据。配置通过测试、应用和 Runtime 状态反馈完成衔接；界面应同时体现配置目标与实际生效状态。Luna 与 Terra 使用同一底层模型 ID 时，必须按档位区分选项，避免选择标准档却被识别为快速档。

真实 Key 保存在部署侧，报告与可提交配置示例只使用占位符。模型服务为外部依赖；作品创新在于群聊运行机制、业务集成和计算策略组织，而非训练或发布新的基础模型。
\source{[S3, S9, S10]；模型 API 字段参考 DeepSeek 官方 Thinking Mode 文档 [W3]，访问日期 2026-09-06。}

\section{Capability 与统一 MCP 设计}
\subsection{业务能力、协议工具与独立插件}
Capability 描述面向任务的业务操作，MCP Tool 描述具体服务暴露的协议调用，AstrBot 插件提供平台接入或独立交互。三者可以关联，但并不一一对应。例如，评课高层业务查询包含多种 operation；校车是本地 Provider；Emoji 与复读使用独立插件入口。

核对时主仓包含 22 份能力定义，其中 21 份配置 MCP mapping，校车通过本地 Provider 装配。定义记录能力身份、输入输出 Schema、Provider、权限、风险和副作用；mapping 记录业务参数如何投影为具体工具调用。原子查询验证的 26 项是测试案例数量，不能替代这一配置数量。

\subsection{能力目录与参数绑定}
能力发现先检查群策略、必要权限及 Provider 条件，再从合法候选中按类别、意图和实体选取业务入口。当前 \code{EntityQueryToolPlanner} 是确定性的单步规划器，将模型提取的实体与明确的查询意图绑定到输入 Schema。

\begin{tabularx}{\textwidth}{@{}L{30mm}Y Y@{}}\toprule
业务约束 & 参数化方式 & 需要保留的区别\\\midrule
课程与教师 & 课程关键词、教师线索、具体条目 ID & 同名课程、不同教师、评论与开课事实\\
活动筛选 & 关键词、类别、时间、数量 & 公开活动、个人记录、报名状态与余位\\
学期与方案 & 学期、年级、专业、方案类型 & 当前学期、历史学期；普通主修与其他方案\\
通知检索 & 内容关键词、条数、详情 ID 或日历范围 & 通知正文、来源字典与类别统计\\
出行查询 & 起点、终点、日期类型、时段 & 行车方向、始发与途经、正常班次与备注\\\bottomrule
\end{tabularx}

以评课为例，本轮自然语言首先形成“评课查询 + 课程关键词”的业务请求，再通过配置映射为 \code{icourse_public_query}。Schema 限制字段、类型、枚举和返回规模。缺少必需信息时应澄清，不能随意构造教师 ID、活动 ID 或方案版本。

\subsection{执行与结果校验}
执行器取得经验证的计划，检查能力仍可执行，调用相应 Provider，并将结果形成 Observation。Observation 同时包含运行状态和业务资料；工具返回空列表仍是一种有效结果，而超时、权限拒绝、返回结构错误属于不同失败。

结果校验负责确保结构、来源和事实字段可以被后续回答使用。最终回答仍可能存在措辞、取舍或语义错误，因此工具 Schema 校验不能被称为“事实正确性全部验证”。项目保留真实问答检查，单独观察是否答对当前问题、是否遗漏来源、是否遵守数量和格式约束。

\subsection{统一客户端与依赖隔离}
MCP 采用 Host、Client、Server 结构，客户端管理与各服务的连接、初始化、能力发现和调用。[W4] 嘟嘟哒在核心侧实现 UnifiedMcpClient，拥有服务目录、工具 allowlist、会话与调用控制。协议中的 capability negotiation 与项目的业务 Capability 是不同层级的概念。

\begin{figure}[H]\centering
\begin{tikzpicture}[node distance=7mm]
\node[box,text width=13cm] (core) {Core：Capability Runtime + UnifiedMcpClient\\业务权限、服务目录、Schema、超时与错误处理};
\node[box,text width=13cm,below=of core] (worker) {独立 Python 解释器：Unified MCP Worker\\mcp 2.0.0；initialize / discover / call / cancel / close};
\node[box,text width=13cm,below=of worker] (services) {独立校园服务进程 / 配置的 MCP 服务\\宿主与 iCourse 使用 mcp 1.29.0；通过 legacy 模式衔接};
\draw[flow](core)--(worker);\draw[flow](worker)--(services);
\end{tikzpicture}
\caption{MCP 版本隔离。worker 负责传输，不接管权限、规划、调度或消息发送。}
\end{figure}

独立 worker 将不同版本依赖隔离在解释器边界，避免为了接入新客户端而整体替换 AstrBot 的依赖。Core 保留计划和授权，worker 只执行已限定的协议操作。这增加一个进程边界，但降低了框架升级对校园服务和宿主插件的影响。

\subsection{扩展一个新服务}
新增服务的顺序为：实现有界输入输出与服务适配；登记 MCP Server 并验证发现和查询；创建业务 Capability definition 与 mapping；在组合中绑定 Provider；最后验证自然语言任务、群配置和失败提示。若只完成 Server 登记，控制台可以检测连接，但 Planner 尚不能使用该服务。

更长的 Skill 将在现有业务能力上组织多步任务。其后续难点包括中间结果引用、跨服务身份对应、总预算、失败后是否继续，以及需要人工确认的动作。当前单步实现提供明确输入输出和回执基础，但尚未完成这一编排层。
\source{[S3, S5, S6, S7]；\nolinkurl{configs/capabilities/}、\nolinkurl{adapters/capability_planner.py} 与 Unified worker README。}

\section{人格、记忆与持续参与设计}
\subsection{Persona 的职责}
Persona 作为独立资产记录语气、句长、词汇、表情预算和场景表达。当前仓库提供 dududa 与 neutral 两套资产，正式组合固定使用嘟嘟哒人格，并以 neutral 作为 Registry 回退资产。资产可扩展不等于已经提供完整的多 OC 切换体验。

回答生成先使用当前问题、工具资料和表达约束产生内容，再由 Composer 组织结构，确定性 Persona Renderer 完成最终处理。人格不应修改课程名称、教师身份、日期、数字、来源或任务目标。其价值在于稳定的信息取舍与表达节奏，而不是固定口号或对每一条回复添加可爱语气。

用户的精确输出要求优先。当前已存验证包含指定短回复、JSON 和多人分工等用例；这些验证只能说明相应用例的行为，不能据此声明所有开放式输出都能严格满足格式。

\subsection{近期上下文与长期记忆的区分}
近期上下文是当前讨论窗口，帮助理解省略、引用和更正；长期 Memory 是跨运行持久化、按范围检索的资料。两者在保留期限、可见性、删除语义和使用目的上不同。

Memory v2 已提供用户画像、群记忆、情节记忆和显式记忆等类型，以及 Repository、检索策略、候选写入筛选、删除、归档、恢复和导出组件。当前生产组合没有把这些组件接入聊天上下文和自动写入。因此，作品可以声称具备记忆基础模块，尚不能声称正式聊天已经持续记住群友或形成可靠关系模型。

\subsection{作用域优先的检索流程}
检索先通过当前 Actor 与 ConversationScope 生成允许的 Selector，限定当前会话、当前群或可用的用户画像。随后处理时间、有效期、可见性和删除状态，得到有限候选集，最后排序。先限制资料范围再计算相似度，可以避免先在全库匹配、再遗漏过滤的问题。

\begin{center}
\begin{tikzpicture}[node distance=5mm]
\node[future,text width=13cm] (a) {检索请求：Actor、ConversationScope、问题、记忆类型};
\node[future,text width=13cm,below=of a] (b) {授权 Selector $\rightarrow$ 有效期 / 可见性 / 删除状态 $\rightarrow$ 有界候选};
\node[future,text width=13cm,below=of b] (c) {近期优先或中文 BM25 $\rightarrow$ Top-K 结果 $\rightarrow$ 后续 Context 接入};
\draw[flow,dashed](a)--(b);\draw[flow,dashed](b)--(c);
\end{tikzpicture}
\end{center}

中文检索实现使用 NFC 归一化、大小写折叠、连续汉字二元组及字母数字词项，无需外部 Embedding 服务。BM25 在同一已授权候选集合上计算：
\[
\operatorname{score}(d,q)=\sum_{t\in q}\operatorname{IDF}(t)
\frac{f(t,d)(k_1+1)}{f(t,d)+k_1(1-b+b|d|/\overline{|d|})},
\]
\[
\operatorname{IDF}(t)=\ln\!\left(1+\frac{N-df(t)+0.5}{df(t)+0.5}\right),
\qquad k_1=1.2,\quad b=0.75.
\]
其中 $N$ 为候选记录数，$df(t)$ 为含词项 $t$ 的记录数，$f(t,d)$ 为词频。当前实现去重查询词项，仅保留正分结果，同分按记忆 ID 排序，便于复现。中文单字、别称与语义等价表达仍可能超出词法匹配能力。

\subsection{写入与生命周期}
记忆写入先形成候选，再由 Write Gate 判断是否允许保存、以何种类型与范围保存。删除后的 tombstone、归档与恢复记录用于防止已删内容在重建过程中重新出现；导出按可授权范围返回。既有测试关注范围隔离、有效期和生命周期行为。

本地合成评测比较无记忆、近期优先和中文 BM25 三种策略，记录 Precision、Recall、MRR 和 nDCG 等指标。但数据卡明确标注没有完成真人相关性评审，因此这些离线结果不能证明真实中文记忆质量或长期体验收益。

\subsection{主动日报与在线学习的阶段}
仓库另有主动任务、调度器、来源存储及摘要预览等组件。它们研究如何在没有新的入站消息时创建有授权、有时间范围的运行，与第 5 节的普通消息触发自动参与属于不同机制。当前提交不宣称已稳定提供面向真实群的自动日报。

Bandit 模块已实现离线决策、反馈与 IPS/SNIPS/DR/ESS 计算，固定四个合成样本用于验证计算和记录行为。它未接入生产 Router，也没有由群友反馈驱动的线上探索。后续只有在建立可归因质量反馈和等价可选端点后，才能评估在线优化是否优于当前静态策略。
\source{[S12, S13, S14]；\nolinkurl{memory/lexical.py}、Memory 数据卡与离线 Bandit 报告。}

\section{数据组织、可靠性与技术难点}
\subsection{状态分层与存储}
系统将源码、公开配置模板、派生数据和私有运行状态分开。源码仓库包含能力定义、映射、人格、依赖锁文件及测试样例；运行目录保存模型凭据、QQ 登录态、账号配置、业务缓存和执行状态。报告不携带这些私有运行材料。

\begin{table}[htbp]\centering
\caption{主要数据与存储责任}
\begin{tabularx}{\textwidth}{@{}L{31mm}L{32mm}Y@{}}\toprule
资料 & 存储或管理方式 & 责任及当前限制\\\midrule
群服务与模型策略 & JSON 配置与应用状态 & 按账号、会话或群读取策略；修改后核对实际生效状态。\\
QQ 登录与历史 & NapCat / AstrBot 私有目录 & 由平台与宿主维护；不进入报告、源码或模板。\\
Web 缓存 & 浏览器 IndexedDB、网关状态 & 缓存已取得的消息及草稿，不成为另一份权威 QQ 数据源。\\
校园缓存 & 服务 SQLite / 研究快照 / 时刻表 & 每个来源保留自身时间与版本；静态结果不得标为实时。\\
单次 Runtime 状态 & 当前组合使用内存 Store & 保存阶段、重复请求与短期结果；不据此声称进程重启后任意阶段自动恢复。\\
入站与投递记录 & Rollout 账本和输出回执 & 记录处理归属与发送结果；未知投递状态需区分，避免盲目重发。\\
Memory v2 & 独立 JSON Repository 等组件 & 已有范围与生命周期实现；未进入正式聊天长期读写。\\\bottomrule
\end{tabularx}
\end{table}

当前组合中的 Runtime 状态 Store 配置为 30 分钟 checkpoint 与重复记录保留窗口，并设置数量上限。这是有界进程内状态管理；与 SQLite 调度器或持久 Rollout 账本承担的职责不同，不能合并表述为完整分布式事务能力。

\subsection{重复消息与输出确认}
平台可能重推消息，用户也可能重试相同请求。核心依照消息与会话身份识别重复执行；输出由独立的 DeliveryRequest 和 Receipt 描述，避免把准备好文本当作已经发送。接口测试覆盖并发重复消息只运行一条模型和工具链，以及重复输出请求复用回执的行为。

实际网络中，超时可能发生在平台已收到消息但本地尚未取得确认之后。系统必须保留失败与结果未知的区别。当前测试与记录可证明特定路径的去重和回执处理，不能据此承诺所有故障下的端到端“恰好一次”投递。

\subsection{会话范围与外部访问}
会话范围贯穿上下文、群配置、能力执行和输出目标。浏览器不能任意指定一个群就读取其完整资料；模型不能用生成文本扩大数据范围。二课凭据只在相应服务进程解析，报告示例使用占位符；开放式文本与工具结果被视为资料，不成为更高权限指令。

当前安全机制包括确定性授权、参数与输出结构验证、受限工具目录、凭据外置和管理入口认证。它们限制程序可以执行的动作，但不构成对所有提示注入、模型幻觉或错误判断的绝对保证。评测应分别检查权限行为、数据处理和回答内容。

\subsection{关键难点与解决方法}
\begin{longtable}{@{}L{26mm}L{63mm}L{60mm}@{}}
\caption{技术难点、实现方法与剩余问题}\\\toprule
难点 & 当前方法 & 效果与剩余问题\\\midrule\endfirsthead
\toprule 难点 & 当前方法 & 效果与剩余问题\\\midrule\endhead
多人短句与更正 & 输入保留身份和引用，读取同群历史，区分当前要求与历史材料，混合结构化感知 & 更正与追问有代表验证；跨话题解缠、全天范围和发言归属仍需提高。\\
异构服务参数 & 业务定义与协议 mapping 分开，确定性绑定实体、日期、类别和数量，校验 Observation & 六类业务已接入；含糊问题仍可能出现先查后澄清或选错操作。\\
群聊打扰 & 群开关、概率、冷却、小时上限和运行占用检查；受授权的规则自适应启用 & 固定讨论与控制逻辑可复现；尚无长期自然群误触发率。\\
计算与时延 & 感知与回答分角色，档位、推理和篇幅正交；暴露模型及耗时 & 真实样本可完成；查询总时延有波动，尚未量化各阶段占比。\\
事实与人格 & 工具事实先组织，Persona 处理表达，格式与可见篇幅最后检查 & 精确格式有回归用例；来源覆盖和开放式回答仍需人工评价。\\
依赖冲突 & MCP v2 放入独立 worker，Core 拥有统一调用，宿主保留原依赖 & 避免整体替换宿主 SDK；进程、会话和超时边界需要联动维护。\\
多入口一致性 & QQ 与 Web 共用组合与 Runtime，替换输出适配器；配置保留应用状态 & 预览可观察实际链路；预览成功仍需与 QQ 投递单独验证。\\\bottomrule
\end{longtable}

这些选择集中解决已出现的故障与交互困难。项目的技术难度主要来自多模块在一次群聊任务中的协同，而不是简单增加工具或模型数量。
\source{[S2, S3, S4, S9, S11, S15]；Runtime Store、投递适配与权限组件。}

\clearpage
\section{WebUI 设计}
\label{sec:webui}
\subsection{设计目标与信息架构}
WebUI 是嘟嘟哒面向管理员的日常工作台。其设计围绕三个连续任务展开：掌握真实群聊与账号状态，为目标群配置智能体行为，检查一次回答的执行过程。界面同时承载 QQ 通信与智能体管理，因此始终保留当前账号、会话及连接状态，让操作对象与结果可以对应。

信息架构分为会话工作区和管理页面。会话工作区让管理员在阅读群聊时直接打开 Agent；管理页面集中处理联系人、请求通知、群服务、人工内测与模型接入。路由保留账号、会话类型和对象标识，刷新页面后仍能定位同一会话。

\begin{table}[H]\centering
\caption{WebUI 的信息架构与主要任务}
\begin{tabularx}{\textwidth}{@{}L{29mm}L{42mm}Y@{}}\toprule
入口 & 管理员任务 & 主要信息与操作\\\midrule
消息工作区 & 阅读和处理当前讨论 & 账号切换、会话列表、消息历史、搜索、输入与回复；可并排展开 Agent。\\
Agent 工作台 & 配置与检查群聊智能体 & 对话、运行、配置三个页签；展示本轮上下文、回答、实际模型和能力调用。\\
联系人与通知 & 定位人员、群与待处理事项 & 群与好友目录、群管理入口、请求及事件；操作绑定具体 QQ 账号。\\
群服务 & 选择群服务初值 & 服务档案、当前状态与操作结果；依赖管理会话及对应控制面后端。\\
人工内测 & 检查既有样本和回答质量 & 浏览已配置的脱敏样本、生成候选并登记人工反馈。\\
API Key 池与设置 & 维护模型接入及工作区 & 三档模型连接、Key 状态、应用到 Runtime；主题、缓存和连接相关设置。\\\bottomrule
\end{tabularx}
\end{table}

\subsection{工作区布局与视觉交互}
桌面工作区采用四列布局：账号与导航栏、会话列表、消息正文、Agent 面板。消息与 Agent 并排，减少管理员在讨论内容和配置结果之间切换的成本；Agent 可折叠，为长消息与资料留出空间。管理页面使用账号栏加主体内容的布局，顶部明确显示执行操作的账号。

\begin{figure}[H]\centering
\begin{tikzpicture}[x=1cm,y=1cm,font=\small]
\draw[draw=rulegray,fill=white,rounded corners=1pt] (0,0) rectangle (15.2,6.2);
\fill[panel] (0.01,5.3) rectangle (15.19,6.19);
\draw[rulegray] (0,5.3)--(15.2,5.3);
\node[anchor=west,font=\small\sffamily\bfseries,text=accent] at (0.25,5.76) {当前账号 / 会话};
\node[anchor=east,text=muted] at (14.95,5.76) {QQ 账号状态\quad Agent 状态};
\fill[panel] (0.01,0.01) rectangle (1.65,5.29);
\draw[rulegray] (1.65,0)--(1.65,5.3);
\draw[rulegray] (5,0)--(5,5.3);
\draw[rulegray] (10.15,0)--(10.15,5.3);
\node[align=center,text width=1.3cm] at (0.825,2.8) {账号\\[4mm]消息\\联系人\\通知\\[4mm]管理\\设置};
\node[anchor=north,align=left,text width=2.8cm] at (3.325,5.02) {\textbf{会话列表}\\[3mm]搜索 / 筛选\\[5mm]选中的群聊\\[3mm]其他群与私聊\\[3mm]最新消息与未读};
\node[anchor=north,align=left,text width=4.55cm] at (7.575,5.02) {\textbf{消息正文}\\[3mm]当前讨论与时间\\[3mm]发言者、引用与消息段\\[3mm]历史分页与搜索};
\draw[rulegray] (5,1.18)--(10.15,1.18);
\node[align=center,text=muted] at (7.575,0.58) {消息输入\quad 提及 / 回复 / 发送};
\node[anchor=north,align=left,text width=4.4cm] at (12.675,5.02) {\textbf{Agent 工作台}\\[2mm]对话\quad 运行\quad 配置\\[3mm]同群近期上下文\\[2mm]候选回答与资料来源\\[2mm]模型、能力与耗时};
\draw[rulegray] (10.15,1.18)--(15.2,1.18);
\node[align=center,text=muted] at (12.675,0.58) {问题输入\quad 运行预览};
\end{tikzpicture}

\caption{桌面工作区的结构示意。图中仅表达布局和信息职责，不使用真实账号或聊天截图。}
\label{fig:webui-layout}
\end{figure}

响应式布局围绕任务切换调整：中等宽度下，Agent 变为可收起的侧边面板；宽度不超过 860 px 时，会话列表、消息和 Agent 以单页切换呈现，底部保留主要导航，其余入口收进“更多”。页面保存选中的会话，避免从 Agent 返回聊天时丢失位置。

界面以低饱和青色标识主要操作，使用浅色背景、细分隔线和有限的状态色组织密集信息，并提供深色主题。在线、警示与失败状态同时给出文字，模型档位和结果类型保持明确标签。键盘焦点可见，搜索和编辑弹窗支持 Esc 关闭、Tab 循环及关闭后恢复焦点；这些细节支持连续管理操作。

\subsection{前后端分工与消息同步}
前端以 Vue 3 和 TypeScript 实现，\code{App.vue} 组合账号栏、会话列表、聊天窗和 Agent 面板；工作区状态由组合式逻辑管理，目录查询与页面数据通过服务适配器进入组件。类型与响应校验集中在数据边界，组件负责呈现、输入和交互反馈。

Node.js Gateway 向浏览器提供同源 HTTP 与 SSE。HTTP 处理工作区读取、历史分页、配置和明确触发的操作；SSE 推送消息及连接变化。网关经认证的 OneBot v11 反向 WebSocket 接收 NapCat 事件，并在服务端执行限定的协议动作。OneBot 连接凭据由网关保管。

\begin{figure}[H]\centering
\begin{tikzpicture}[node distance=8mm and 8mm]
\node[box,text width=13.5cm] (ui) {Vue 浏览器：会话工作区、Agent 工作台与管理页面};
\node[box,text width=13.5cm,below=of ui] (gateway) {Node.js Gateway\\工作区快照 / SSE 事件 / 历史读取 / 配置与预览请求};
\node[box,text width=6.1cm,below=11mm of gateway.south west,anchor=north west] (qq) {NapCat / OneBot v11\\账号、消息与 QQ 操作};
\node[box,text width=6.1cm,below=11mm of gateway.south east,anchor=north east] (runtime) {AstrBot / Dududa Runtime\\感知、模型、校园工具与回答};
\draw[<->,>=Stealth,draw=accent,thick] (ui)--node[right,labeltext]{同源 HTTP / SSE}(gateway);
\draw[<->,>=Stealth,draw=accent,thick] (gateway.south -| qq.north)--node[left,labeltext]{反向 WebSocket}(qq.north);
\draw[<->,>=Stealth,draw=accent,thick] (gateway.south -| runtime.north)--node[right,labeltext]{Runtime HTTP API}(runtime.north);
\end{tikzpicture}

\caption{WebUI 的数据与调用关系。QQ 同步和 Agent 预览共用网关入口，分别连接各自的后端。}
\label{fig:webui-architecture}
\end{figure}

首屏先读取已有工作区快照，联系人与历史在后台继续加载；慢联系人或慢 Agent 请求不会阻止会话框架显示。SSE 重连根据事件位置补发可重放事件，超出保存范围时刷新工作区。历史分页与账号、会话绑定，消息按规范化标识合并。

浏览器使用 Dexie/IndexedDB 缓存已取得的消息、会话、历史覆盖范围和草稿，主键包含账号与会话标识。搜索明确标注当前浏览器缓存范围；未取得的离线消息不被描述为已经同步。缓存清理作用于本地副本，QQ 数据仍以 NapCat 返回结果为依据。

\subsection{Agent 对话、运行与配置}
Agent 面板将使用过程拆成三个页签：\textbf{对话}接收管理员请求并呈现回答与资料；\textbf{运行}展示实际模型、推理与篇幅配置、使用的能力、耗时、上下文覆盖和错误原因；\textbf{配置}维护本群的行为偏好、插件开关与参与策略。一次请求的运行信息与回答保持关联，管理员可以从结果直接追查执行过程。

在当前正式 Runtime 预览路径中，请求依次经过以下步骤：
\begin{enumerate}
\item 浏览器提交当前 \code{accountId}、\code{conversationId} 和问题。
\item 网关校验输入与会话范围，从服务端读取同群近期历史；读取失败时直接说明原因。
\item 网关使用服务端凭据调用 AstrBot 的 Runtime 预览接口，由既有核心完成感知、路由、必要查询和回答组织。
\item Runtime 返回回答、模型、能力、耗时和上下文用量，网关校验后交给页面呈现。
\item 管理员在“运行”页查看此次实际执行情况，再决定是否调整配置或继续提问。
\end{enumerate}

预览模式可执行校园工具查询，但结果留在 Web 中，返回的 QQ 输出调用与长期记忆写入均为 0。QQ 手动消息发送属于聊天窗的独立操作。人工内测页则面向已配置样本和人工评价，其结果与正式 Runtime 预览分别记录。

配置以账号和会话为作用域，模型档位、推理强度、回答长度、回复强度、上下文预算和群聊风格分别设置。页面区分管理员偏好、允许范围与锁定状态，并展示本轮有效选择；“上下文长度”指历史读取预算，不能解释成模型服务商的最大上下文窗口。最终使用值以 Runtime 返回结果为准。

\subsection{模型、工具与服务管理}
模型管理按 Luna、Terra、Sol 三个档位组织连接和 Key 池，显示真实模型 ID、接入地址、推理参数、超时、输出额度及 Key 可用状态。相同底层模型可以服务不同档位，选项使用档位标识区分，避免按模型名称合并后选择错位。

配置编辑、保存和运行生效有各自的反馈。模型连接保存后，页面读取 Runtime 配置状态；管理员执行“应用到 Runtime”后，等待服务端确认，再用一次请求观察实际模型。这样可以定位“内容已保存，但运行侧仍使用旧配置”的问题。

\begin{table}[H]\centering
\caption{配置对象与生效反馈}
\begin{tabularx}{\textwidth}{@{}L{28mm}Y Y@{}}\toprule
配置对象 & 页面操作 & 需要观察的结果\\\midrule
群行为与插件 & 选定会话，设置参与方式、偏好和插件范围并保存 & 当前群配置、可用能力、本轮实际选择及启用原因。\\
模型与 Key 池 & 编辑档位连接，维护 Key，检查后应用到 Runtime & 待应用或已应用状态、运行侧是否就绪、回答实际使用的模型。\\
MCP 与扩展插件 & 查看目录，检查连接，按部署能力安装或配置 & 服务连接结果、工具目录与错误说明；是否进入 Planner 另由能力映射决定。\\
群服务档案 & 在控制面选择档案，预览并提交配置 & 当前服务分配、命令结果与后端可用状态。\\\bottomrule
\end{tabularx}
\end{table}

MCP 管理界面提供服务检查与接入，插件管理提供安装和运行状态入口。服务出现在目录中表示已经登记，能否参与智能体回答还取决于当前能力映射、授权与运行配置。群服务页和历史内测页若缺少后端或样本，会展示配置说明及重新检测入口。

\subsection{连接状态、异常与可用性}
WebUI 分别呈现网关连接、QQ 账号在线状态、Agent 预览可用性和生产自动参与状态。网关可访问时，QQ 仍可能需要重新登录；预览接口可用时，自动参与也可能尚未开启。这些状态决定管理员下一步应检查账号、模型还是群策略。

\begin{table}[H]\centering
\caption{试用问题如何转化为交互设计}
\begin{tabularx}{\textwidth}{@{}L{34mm}Y@{}}\toprule
实际问题 & 对应交互与实现处理\\\midrule
首屏等待联系人刷新 & 先显示工作区快照，后台刷新目录；失败时保留已读取内容并提供重试。\\
QQ 离线仍显示连接正常 & 显示在线账号数与账号总数，离线时给出登录指引和重新检测。\\
Agent 预览失败原因不明 & 区分历史读取与 Runtime 调用失败；在对话及运行页展示原因和重试。\\
目录失败被误认为空列表 & 分开表示未读取、读取失败与真实空结果；旧数据标记为上次读取结果。\\
搜索范围和焦点不清楚 & 标注浏览器缓存范围，区分未搜索、加载和无匹配；完善键盘与焦点恢复。\\\bottomrule
\end{tabularx}
\end{table}

默认工作台入口绑定本机回环地址；对外部署由反向代理提供登录保护，群服务管理另外依赖操作员会话与后端授权。页面仅展示管理员完成当前操作所需的信息，模型 Key、OneBot 凭据和私有文件路径由服务端处理。

\subsection{典型操作流程与实现依据}
以启用选课辅助为例：管理员选择 QQ 账号和目标群，确认账号及 Runtime 状态，开启已授权的评课服务与合适的参与方式；随后在 Agent 对话中输入代表问题，在运行页检查实际模型、评课能力、上下文与资料来源；若回答篇幅或参与节奏不合适，再返回配置页调整。工作区把“看到讨论—设置服务—检查回答”组织为连贯操作，支持校园场景中的持续维护。

本章以当前页面、路由、网关与 Runtime 适配器源码为主要依据，并结合开放式试用修复记录。旧版 Web README 中有关早期候选模式和服务数量的描述仅作为历史材料；本章不以其替代当前调用路径。Web 的已有回归结果见第\ref{sec:evaluation}节。
\source{[S16, S18]；前端组件、路由、缓存、网关与 Runtime 适配器，以及开放式试用修复记录。完整路径见附录。}

\clearpage
\section{部署与交付方案}
\label{sec:deployment}
\subsection{部署拓扑}
\begin{figure}[H]\centering
\begin{tikzpicture}[node distance=7mm and 7mm]
\node[box,text width=6cm] (web) {Web 容器\\Vue 静态资源 + Node Gateway\\宿主默认入口 5173};
\node[box,text width=6cm,right=of web] (nc) {NapCat 容器\\QQ 登录与 OneBot\\管理入口 6099};
\node[box,text width=6cm,below=of web] (mc) {MCP Console 容器\\服务目录、连接诊断、查询};
\node[box,text width=6cm,right=of mc] (bot) {AstrBot 容器\\Dududa Core + 模型适配\\管理入口 6185};
\draw[flow](nc)--(web);\draw[flow](nc)--(bot);
\draw[flow](web.south east)--(bot.north west);
\draw[flow](web)--(mc);
\node[box,text width=13cm,below=of bot.south west,anchor=north] (vol) {运行卷：QQ 登录态、AstrBot 配置、群策略、校园缓存、Key Store};
\end{tikzpicture}
\caption{默认 Compose 的四个服务。模型供应商和校园数据源位于外部；图中端口为本地部署入口。}
\end{figure}

AstrBot 容器装配智能体核心与插件，MCP worker 由独立解释器执行；校园服务源码与配置通过镜像或只读挂载提供，业务缓存与运行状态写入独立目录。多 QQ 账号需要各自 NapCat 登录实例和独立数据卷，可反连同一 Web 网关。

\subsection{交付组成与复现入口}
程序采用源码、核心 wheel、Web 构建产物与 Docker Compose 交付，随包提供依赖锁文件和运行说明。开发环境使用 Python 3.12、uv，以及 Node.js 22.18.0/npm 10.9.3。精确依赖由 \code{uv.lock} 和 Web 锁文件确定。

\begin{Code}
packages/dududa-agent/   Core: domain, ports, runtime and modules
apps/astrbot-plugins/    AstrBot adapters and owned plugins
apps/web/               Vue frontend and Node gateway
services/mcp/           Campus servers, console and worker
configs/                Capabilities, mappings and persona assets
deploy/                 Compose, Dockerfiles and examples
ops/                    Installation, operation and evaluation
tests/ and evals/       Behavioral checks and datasets
docs/                   Design, operation and submission reports
\end{Code}

无需 QQ 登录或模型 Key 的离线验证入口为：
\begin{Code}
uv sync --locked --python 3.12.13
uv run --locked python ops/cli/run_dududa_100_message_benchmark.py \
  --json-output /tmp/dududa-100-runtime.json \
  --report-output /tmp/dududa-100-runtime.md
\end{Code}

完整安装时，部署者按照程序使用说明配置环境文件、运行目录、模型服务与必要的校园认证，执行初始化与启动流程，再完成 NapCat QQ 登录、AstrBot Provider 和目标群策略设置。报告提供的是结构与复现入口，具体部署步骤以随程序的 \code{submission-program.md} 为准。

升级时应保留现有运行卷与旧镜像，并依据受影响服务进行替换；运行中的账号、凭据和群授权不属于报告生成过程。模板与报告的编译方法在本报告源码目录另行提供。
\source{[S8, S9, S10, S16]；Web README、Compose 与程序使用说明。}

\section{验证结果与完成度}
\label{sec:evaluation}
\subsection{评测方法与证据分层}
本地材料包含不同阶段、模型和执行模式的测试，不能合并为一个总体通过率。离线固定消息用于验证程序路径；真实模型预览用于检查语义和回答；原子 MCP 测试检查具体服务；模拟入站用于验证事件转换与参与；QQ 投递则需独立查看输出记录。

本报告引用已存资料时保留原测试日期和限制，并在报告编制过程中重新执行一轮不联网的 100 条固定消息流程。未为撰写报告重新向 QQ 群发送消息，也未调用付费模型生成新测试答案。验证中涉及的真实查询结果均来自既有记录。

\begin{longtable}{@{}L{26mm}L{61mm}L{62mm}@{}}
\caption{主要证据及其适用范围}\\\toprule
证据 & 输入与执行范围 & 能够支持的结论\\\midrule\endfirsthead
\toprule 证据 & 输入与执行范围 & 能够支持的结论\\\midrule\endhead
100 条固定消息 & 脚本模型、校园 fixture、内存输出，经过生产组合；真实 QQ 发送为 0 & 验证路径、参数投影、调用次数、结果状态及去重；不衡量真实模型回答准确率。\\
26 项 MCP 查询 & 录制环境真实服务的原子调用；包含可选服务案例 & 26 个选定案例通过，证明相应服务可调用；不代表全部自然语言任务。\\
5 项自然语言问答 & 真实 DeepSeek Flash、真实业务资料、Runtime 预览 & 代表问题调用预期能力并得到已核对的结果；样本较少且非随机抽样。\\
四人 20 条讨论 & 固定台词、模拟发言者、正式参与路径、真实模型与评课服务 & 代表场景能从关闭评课到授权启用，再查询并回应群成员需求。\\
原生入站排练 & OneBot 转换与正式接收路径；记录 1 次 output call 和 canary 完成 & 已存记录支持一次回复路径完成；模拟成员输入不能称为四个真实用户账号完整发言。\\
广覆盖群聊预览 & 93 个真实模型预览请求，另 7 项隔离测试 & 暴露语义、格式、来源和上下文问题；保留部分通过与失败结果。\\
Web 与组件回归 & 指定版本的自动化和浏览器场景 & 支持相应回归行为；不能证明所有线上服务配置、所有真实群或所有设备。\\\bottomrule
\end{longtable}

\subsection{离线 100 条固定消息流程}
既有 2026 年 9 月 6 日录制结果为 94 条 Runtime 完成响应，6 条由入口规则保留为兼容路径处理。执行统计为 187 次脚本模型调用、53 次 MCP fixture 调用和 11 次本地校车能力调用；真实 QQ 发送、Memory 写入均为 0。

动作、结果、投递预期、工具名、能力步骤和重复重放检查均无不匹配。每条业务计划不超过一步。这里的“兼容路径”是测试中记录的入口处置，不表示旧 1.0 智能体在生产中继续处理这些请求。

编制报告时按第\ref{sec:deployment}节命令重新运行同一 fixture，结果见随源码提供的离线基准摘要与检查记录。这轮复测用于校验报告中的流程统计与当前工作区一致，不增加真实模型质量样本。

\subsection{真实模型代表任务}
以下数据来自已存 \code{manual-answers.json}，运行于 2026 年 9 月 6 日录制验证。使用 Flash 标准档、关闭额外 thinking，实际读取业务资料。耗时是单次观测，包含整体请求过程，不是平均时延、P95 或性能承诺。

\begin{table}[htbp]\centering
\caption{五项真实查询的已存结果}
\begin{tabularx}{\textwidth}{@{}L{37mm}Y L{19mm}@{}}\toprule
任务 & 结果要点 & 单次耗时\\\midrule
2026 秋数学分析开课 & 返回 A1/B1 开课记录、教师和学期 & 10.92 秒\\
2026 级计算机普通主修 & 返回对应方案及总学分 164；绑定方案口径 & 12.06 秒\\
最近三条校园通知 & 返回三条通知、日期与来源链接 & 12.02 秒\\
工作日东区至西区上午校车 & 返回相应方向与时段班次 & 17.04 秒\\
二课英语角活动 & 返回 English Corner 活动、时间和报名状态 & 29.07 秒\\\bottomrule
\end{tabularx}
\end{table}

上述 5 项记录均通过对应检查。由于样本来自演示任务，无法代表所有校园问题，也不能用这些耗时比较 Flash 与 Pro 的优劣。远端内容会更新，复现时应检查字段、来源与能力选择，而非要求模型逐字输出相同文本。

\subsection{自适应启用与入站排练}
固定案例由四个模拟发言者逐步讨论数学分析选课，包含讲解、证明训练、课业负担和候选筛选等不同需求。前 19 条没有启用评课，第 20 条形成明确请求后，记录启用 \code{icourse.read} 并查询 \code{icourse.public-query.v2}。

原生入站排练记录包含 20 次状态变化、1 次自适应启用、1 份回答、1 次输出调用和 0 次 Memory 写入，最终处置为 \code{canary_completed}。台词注入接收端，不能据此声称四位成员的账号都经 QQ 网络真实发送了 20 条消息。该案例验证的是既定场景下的事件转换、参与、启用、查询与回复路径，长期自然群的误启用率尚未建立。

\subsection{广覆盖预览中暴露的问题}
2026 年 9 月 4—5 日的群聊 100 题材料包含 93 次真实模型合成历史预览和 7 项隔离故障或策略测试。真实预览终态为 89 次 response、4 次 failed；语义审查则为 60 通过、25 部分、8 失败。加上隔离测试登记，合计为 63 通过、29 部分、8 失败。不同层级登记不能写成生产端到端通过率。

问题涉及引用理解、精确格式、查询前澄清、来源缺失、摘要时间范围与发言归属，以及生成失败时正文为空等。后续录制和专项回归修复了其中若干路径，但没有与这批广覆盖预览同范围、同方法的全量重测证据。因此，新版仍保留这组结果，不以“5 项通过”覆盖较早的广泛缺陷。

另有 8 月的评课 75 题审校、二课 75 条消息模拟及 9 月 3 日固定消息报告。它们分别使用过不同模型、脚本语义或旁路审校，适合作为功能演进与回归资料，不能混算成同一 DeepSeek 评测集。

\subsection{组件与界面验证}
录制操作单记载：26/26 MCP 原子查询、5/5 自然语言问答、多个 Python 回归集合，以及 Web 114 项前端和 121 项服务端测试通过。独立 Web 试用修复报告还记录了桌面、手机与键盘场景。不同批次之间可能包含重复测试，因此本文不将这些数量累加为独立测试总数。

Memory 检索与 Bandit 评测均为合成离线数据。它们支持组件行为与计算正确性，不支持长期记忆效用、真实社交反馈或在线路由收益。当前还没有经过设计的用户研究来量化学习效率提升或管理成本下降。

\subsection{当前完成度与后续验证重点}
\begin{tabularx}{\textwidth}{@{}L{34mm}Y@{}}\toprule
阶段 & 能力与验证重点\\\midrule
当前运行能力 & QQ/预览入口、近期上下文、混合感知、确定性社交决策、按群自动参与、授权自适应启用、六类校园查询、模型档位与 Web 管理。\\
独立组件或离线成果 & Memory v2、主动调度与摘要模块、离线 Bandit、扩展 Persona 与可选服务资产。需要进一步接入和真实体验验证。\\
仍需建设 & 多步 Skill、长期群友印象、多 OC 切换、自然群参与质量、摘要时间与归属、查询延迟分解、当前广覆盖真实模型复测。\\\bottomrule
\end{tabularx}
\source{[S10, S11, S12, S13, S16]。数字直接取自已存 JSON、数据卡与结果报告；本轮离线复测单独留存。}

\section{创新点、资料整合与演进路线}
\label{sec:synthesis}
\subsection{作品的创新点}
嘟嘟哒的核心设计主线是\textbf{自适应}：以当前讨论与任务为输入，在管理员设定的范围内选择参与方式、业务插件和计算路由，并通过 WebUI 呈现实际结果。自适应分别落实到能力选择、计算选择和交互选择，形成可解释的群聊服务机制。

\textbf{自适应插件选择。}系统综合讨论主题、多人参与、需求持续性与当前请求，在管理员预先授权的只读服务中选择并启用匹配插件。以选课讨论为例，群成员逐步提出讲解方式、证明训练和课业负担等条件，当前请求达到触发条件后，系统启用评课插件并进入课程资料查询。能力选择因此与连续讨论关联，减少用户记忆命令和手动切换服务的负担；启用原因与实际能力调用可在运行记录中核对。实现机制见第 5.4、7 节。

\textbf{任务驱动的自适应路由。}感知结果提供意图与复杂度信号，档位策略结合置信度、可用模型和预算选择本轮计算配置，路由器再校验角色、能力与输出容量。简单交流、校园查询和复杂比较可以使用不同档位；推理强度与可见篇幅分别控制，避免仅按答案长短分配计算资源。自适应体现在运行选择随任务改变，选择依据可记录并在 WebUI 中查看。当前策略由规则和配置定义，其工程目标是匹配任务所需的计算投入；尚未完成同条件成本与效果对照实验。机制见第 6 节。

\textbf{自适应群聊参与与表达。}系统根据提及、追问、持续讨论和群策略决定回复、查询、澄清或静默，再结合任务要求组织回答篇幅与人格表达。参与选择与冷却、频率限制分别执行，既利用多人交流中的情境信号，也保持管理员对参与节奏的控制。工具事实与 Persona 分层，使严谨解释和日常表达能够在同一运行路径中按需组织。机制见第 5、8 节。

\textbf{WebUI 工作台设计。}工作台将真实 QQ 会话与 Agent 对话、运行、配置置于统一界面，并在桌面并排布局与手机单页切换中保留当前账号和会话。管理员按群设置偏好、允许范围与锁定项，运行后查看实际模型、插件、上下文覆盖、耗时和选择原因，再调整策略。QQ 与 Web 预览共享核心，构成“配置—预览—观察—调整”的连续操作，使自适应机制既能被配置，也能被检查。详细设计见第\ref{sec:webui}节。

这些创新属于校园群聊场景下的系统设计与交互机制创新。当前自适应由语义感知、确定性规则、管理员配置与运行状态共同实现；在线 Bandit、长期记忆学习和策略更新另有独立研究与接入计划。

\subsection{本地资料汇总范围}
本次扫描主仓、主仓阶段记录、两套历史原型/插件目录、原提交材料、Web 试用报告、嘟嘟哒文档站、相关部署记录、录制基准，以及培养方案数据说明。共登记 441 份当前主仓及独立来源文档；加上工作树与发布副本为 8,944 份物理文件，按逻辑路径归并为 446 个条目。数值反映文档整理范围，不是独立研究或测试数量。

完整清单见同目录的“本地文档汇总.md”，每项记录标题、主要议题、采用方式与副本数量；\code{source-inventory.json} 保留各副本位置。报告主体按主题综合材料，避免照搬按时间累积的开发日志。私有运行配置、聊天数据库、凭据、第三方依赖文档和批量培养方案数据条目不作为技术报告输入。

\begin{longtable}{@{}L{29mm}L{62mm}L{58mm}@{}}
\caption{文档体系的整合方式}\\\toprule
资料类型 & 主要贡献 & 在本报告中的使用\\\midrule\endfirsthead
\toprule 资料类型 & 主要贡献 & 在本报告中的使用\\\midrule\endhead
设计文档与 ADR & 核心依赖、消息与状态、模型、能力、记忆、人格、安全、部署边界 & 第 3—\ref{sec:deployment}节；以当前组合根验证已接入内容。\\
研究报告 & 语义感知、群体情境、主动参与、记忆检索、离线策略与插件化方向 & 第 5—9、\ref{sec:synthesis}节；保留为设计动机或后续方向，不转成运行能力。\\
重构与阶段记录 & 从早期插件组合到统一智能体核心的迁移、接口收敛、真实纵向接入 & 解释历史路径、类名与职责变化；旧的“当前状态”按日期使用。\\
服务与插件 README & 服务输入输出、认证和缓存、依赖、显式命令与默认开关 & 第 4、7、\ref{sec:webui}—\ref{sec:deployment}节；区别独立插件、可选 MCP 与 Planner 能力。\\
评测、录制与试用 & 固定消息、真实查询、回答审校、模拟入站、UI 问题及修复 & 第\ref{sec:evaluation}节；分别保留执行模式、日期、失败项和样本限制。\\
文档站与历史原型 & 早期模块设计、13 阶段流水线、8 角色路由、开发与部署经验 & 对照采用接口分工思想；不移植旧原型的功能和测试数字。\\
原提交 PDF / Word & 作品定位、演示配套信息、旧报告与简介 & 统一成熟技术语体，补充架构、API、验证和完成度说明。\\\bottomrule
\end{longtable}

\subsection{资料中的主要差异及处理}
\textbf{旧状态与新实现。}文档站的架构页面固定在早期提交，\code{current-state.md} 也包含重构前基线；其中“尚未实现”或旧插件清单不能覆盖当前组合根。新版按现有源码描述运行路径，历史状态用于解释演进。

\textbf{原型与主仓。}历史 \code{dududa20-prototype} 和 \code{dududa20-plugin} 记录过多步规划、八角色路由、个人信息及记忆等功能，并有自己的测试数字。当前主仓采用不同的实现、组合和阶段划分，新报告不把它们视为同一发布物。

\textbf{设计目标与组件接入。}完整 Memory、主动日报、多 OC 和 Skill 组合在多份资料中反复出现。源码存在、离线验证、配置登记、正式接入与长期效果是不同证据，当前状态已在第\ref{sec:evaluation}节统一。

\textbf{模型配置与测试口径。}Luna/Terra/Sol 是档位别名，旧材料中还使用过其他模型；本次演示记录使用 DeepSeek。常规推理强度与录制时关闭 thinking 的配置分别注明，不把测试数量或样本时延外推到全部模型。

\subsection{下一阶段的技术路线}
近期优先改进来源覆盖、查询前澄清、近期摘要的时间与归属，分解等待时间，并对现有广覆盖题集完成同范围的真实模型复测。参与机制需要在获准的自然群窗口上评估时机、回应对象、误触发和重复建议，而不只评估是否成功生成文字。

随后接入 Memory 的显式写入、检索、删除与反馈闭环，在真实但经授权的数据上比较无记忆、近期优先和 BM25。多 OC 与群友印象应在有明确资料来源和修改方式的基础上建设。

较长远的扩展是以现有业务能力为基础的多步 Skill，以及同角色、同任务范围内的路由优化。二者均应从具体任务和可测收益出发，保持当前业务的可维护性，并保留用户与管理员对群聊参与的控制。
\source{[S1—S18]；文档全集与副本索引随 LaTeX 源码交付。}

\clearpage\appendix
\section{本地资料与源码索引}
以下路径以主仓根目录为基准；录制与历史目录位于主仓旁。完整本地条目、主题及副本位置见配套汇总文件。源码链接以提交时的主仓为准，仓库访问权限由项目维护者管理。

\begingroup\small\renewcommand{\arraystretch}{1.15}
\begin{longtable}{@{}L{9mm}L{44mm}L{96mm}@{}}
\caption{正文来源编号}\\\toprule
编号 & 主题 & 主要文件或目录\\\midrule\endfirsthead
\toprule 编号 & 主题 & 主要文件或目录\\\midrule\endhead
S1 & 作品与原材料 & \nolinkurl{README.md}；\nolinkurl{docs/design/dududa-2.0-design-report.md}；\nolinkurl{docs/design/dududa-2.0-work-introduction.md}；原提交 PDF/Word。\\
S2 & 核心架构与接口 & \nolinkurl{docs/design/runtime.md}；\nolinkurl{packages/dududa-agent/src/dududa/domain/}、\nolinkurl{ports/}、\nolinkurl{runtime/}。\\
S3 & 当前运行组合 & \nolinkurl{apps/astrbot-plugins/astrbot_plugin_dududa_core/composition.py}；\nolinkurl{rollout_bridge.py}；\nolinkurl{web_runtime.py}。\\
S4 & 感知与参与 & \nolinkurl{docs/design/perception-and-social.md}；\nolinkurl{adapters/proactive_talk.py}、\nolinkurl{adapters/adaptive_plugins.py}；核心 \nolinkurl{perception/}。\\
S5 & 能力与工具 & \nolinkurl{docs/design/capability-and-mcp.md}；\nolinkurl{configs/capabilities/}；\nolinkurl{adapters/capability_planner.py}。\\
S6 & 校园服务 & \nolinkurl{services/mcp/} 各服务 README；校车插件 README 和时刻表；培养方案数据 README。\\
S7 & 扩展与整合 & \nolinkurl{docs/integrations/pr10-selective-integration.md}；各自有插件 README；\nolinkurl{THIRD_PARTY_NOTICES.md}。\\
S8 & 安装与部署 & \nolinkurl{docs/operations/submission-program.md}、\nolinkurl{deployment.md}、\nolinkurl{upgrade-and-rollback.md}；\nolinkurl{deploy/compose/compose.yml}。\\
S9 & 模型与 Key 池 & \nolinkurl{docs/design/model-routing.md}；\nolinkurl{docs/development/api-key-pool-runtime-sync.md}；\nolinkurl{adapters/model.py}、\nolinkurl{deepseek_config.py}。\\
S10 & 录制验证 & \nolinkurl{docs/operations/recording-runbook.md}；\nolinkurl{../dududa-recording-state/reports/} 中 manual-answers、manual-mcp、adaptive、native-result、benchmark-100 记录。\\
S11 & 广覆盖评测 & \nolinkurl{docs/operations/group-chat-100-question-results.md}；\nolinkurl{docs/refactor/} 内评课 75 题、二课 75 题、自然语言 MCP 及固定消息报告。\\
S12 & Memory 与 Persona & \nolinkurl{docs/design/memory.md}、\nolinkurl{persona.md}；\nolinkurl{docs/research/memory-evaluation.md}；\nolinkurl{evals/memory-retrieval/v1/DATA_CARD.md}。\\
S13 & 离线 Bandit & \nolinkurl{docs/refactor/s20-offline-bandit-report.md}；\nolinkurl{evals/bandit-offline-v1/}；\nolinkurl{packages/dududa-agent/src/dududa/bandit/}。\\
S14 & 主动调度与研究 & \nolinkurl{docs/design/proactive-messaging.md}、\nolinkurl{online-learning.md}；\nolinkurl{docs/research/} 内群体情境、感知、主动参与及核心模块研究。\\
S15 & 安全与阶段审查 & \nolinkurl{docs/design/security.md}；\nolinkurl{docs/adr/}；\nolinkurl{docs/reviews/}；\nolinkurl{.TreeWork/} 阶段记录。\\
S16 & Web 与运行修复 & \nolinkurl{apps/web/README.md}；\nolinkurl{docs/reviews/webui-exploratory-fixes-2026-09-05.md}；\nolinkurl{docs/operations/bot-stable-release.md}；独立试用报告。\\
S17 & 文档站与旧原型 & \nolinkurl{../mmdustc.top/apps/dududa-docs/}；\nolinkurl{../dududa20-prototype/}、\nolinkurl{../dududa20-plugin/}；主仓重构与迁移文档。\\
S18 & WebUI 设计与实现 & \nolinkurl{apps/web/src/App.vue}、\nolinkurl{router.ts}、\nolinkurl{components/AgentConsole.vue}、\nolinkurl{views/ApiKeyPoolsView.vue}、\nolinkurl{services/database.ts}；\nolinkurl{apps/web/server/app.ts}、\nolinkurl{dududa-runtime.ts}、\nolinkurl{internal-test.ts}；\nolinkurl{docs/design/bot-control-plane.md}。\\\bottomrule
\end{longtable}
\endgroup

\section{外部资料与排版来源}
本次联网核对以下官方项目或文档。外部资料用于模板选择与协议字段解释，项目实现和验证结果仍以本地代码及记录为依据。访问日期均为 2026 年 9 月 6 日。

\begin{description}[leftmargin=3em,style=nextline]
\item[W1] ElegantLaTeX. ElegantPaper：中文及英文工作论文模板，v0.12（2026-02-27）。\url{https://github.com/ElegantLaTeX/ElegantPaper}。
\item[W2] CTeX 项目. 中文 LaTeX 类与宏包说明。\url{https://ctan.org/pkg/ctex}。
\item[W3] DeepSeek. Thinking Mode：thinking 开关、reasoning\_effort 与模型参数。\url{https://api-docs.deepseek.com/guides/thinking_mode/}。
\item[W4] Model Context Protocol. Architecture，2025-11-25 规范。\url{https://modelcontextprotocol.io/specification/2025-11-25/architecture}。
\end{description}

\subsection*{模板采用与修改说明}
本报告实际加载 ElegantPaper 上游类文件；原始文件随源码放在 \code{latex/vendor/}，遵循 LPPL 1.3c 或更新版本，许可证一并提供。中文排版采用 CTeX 与 Fandol 字体，图表由 LaTeX/TikZ 生成。项目在独立 \code{preamble.tex} 中配置 A4 单栏、页边距、标题层级、页眉页脚、表格及深灰/藏蓝配色，没有修改上游类文件。

模板对照与选用理由另见 LaTeX 源码 README。图示优先表达组件责任与运行关系。
\endinput

\end{document}
